\section{背景}
\label{sec:background}

\subsection{告警分诊}
告警分诊紧随检测之后：对每条告警，分析师判断实际发生了什么、是否需要响应~\citep{nodoze}。攻击调查从一个关注点出发重建完整的攻击过程~\citep{kairos,clouseau}；分诊则以单一结论结束，即告警的起因以及支撑它的证据。分析师通过查询环境得出结论。不同查询的成本不同，区分候选解释的能力也不同。有些查询便宜，并且对某一种解释有决定性；另一些，例如威胁情报查询，既慢，又对多种攻击返回同样的答案。因此查询的顺序决定调查多快结束，分析师还必须判断证据何时已经足够。

设想一条告警报告 Web 层认证失败和 HTTP 错误激增。它看起来像撞库（credential stuffing），但同样的症状也可能来自 SQL 注入探测、经授权的漏洞扫描、被误判为攻击的健康检查，或暴露在公网上的管理面板。对头部来源做威胁情报查询，可能会说它是恶意的——而这个答案是其中三种起因共有的。读取管理面板的配置，只需一次便宜的查询，就能立刻判定最后一种解释；但只有始终把这种解释放在视野里的调查者才会去运行它。

\subsection{本地 LLM 智能体}
LLM 智能体把语言模型与工具结合在一起。每一步，模型读取任务和已有观测，然后调用一个工具或结束任务~\citep{react}。在此基础上的扩展加入了语言化的自我反思~\citep{reflexion}，或学习何时调用工具~\citep{toolformer}。在安全领域，这类智能体已被用于渗透测试~\citep{pentestgpt}、夺旗赛（CTF）题目~\citep{cybench}，以及基于安全日志的调查~\citep{excytin,clouseau}，大多使用大型托管模型。Qwen2.5 和 Llama~3.1 等开源权重模型则可以在组织自己的硬件上运行；经过 4 比特权重量化，最大 72B 参数的模型可以放进单块 GPU。本地部署让遥测数据留在本地，代价是模型能力变弱，而在智能体循环中，每一步的决策都依赖这个更弱的模型。
